% =====================================================================
\section{Introduction}
\label{sec:intro}

Visual correspondence (dense optical flow, dense geometric matching, and sparse
semantic keypoint matching) has come to rely heavily on synthetic training data, as
ground-truth correspondence is scarce for real images
\cite{mayer2016FlyingThingsDataset,sintel,greff2021kubric,pointodyssey}. Most
progress targets \emph{how} to render this data (fidelity, lighting, speed)
\cite{greff2021kubric,pointodyssey}, building on the assumption that closing the
appearance gap, through photorealism and scenes filled with realistic, target-like
objects, drives transfer. Yet non-photorealistic data often transfers well
\cite{Baradad2021LookingAtNoise,Kataoka2020Pre-trainingImages,Anderson2022ImprovingFractal,white2009mandelbulb}. As rendering becomes more efficient
and the appearance gap narrows, the open question may no longer be how to render but
whether \emph{motion} matters, and if so, how much.

We study the effects of motion on correspondence across three tasks (dense optical flow, dense geometric matching and sparse semantic keypoint matching) by training four architectures on 11
synthetic and real training sources and measuring transfer to 9 target benchmarks
(Table~\ref{tab:study}). Across the study, how well a source's motion \emph{covers}
the target's motion is a good predictor of transfer, while appearance similarity is not. The
signal is directional. We find coverage discrepancy, how much of the target's motion the
source fails to reproduce, carries it. Off-target mass, the motion a source produces
that the target never uses, does not. Symmetric distances often dilute the signal by
averaging these two distances together (see Fig.~\ref{fig:concept}). Because coverage discrepancy
needs only motion geometry, with no rendering or model training, so it is efficient
enough to act on.

\begin{figure}[tbp]
\centering
\includegraphics[width=\linewidth]{figures/concept_directional_vs_symmetric.png}
\caption{\textbf{Direction matters.} An illustrative two-mode example on toy
distributions, not measured data. A symmetric distance such as the Chamfer
averages the two directions and reads both scenarios below as the same distance,
while the directed distances swap between them and reveal which way the source
and target fail to align.
\textbf{(a)}~\emph{Under-coverage}: the target places mass on a mode the source
never renders, so the coverage discrepancy ($\dbt$) rises while off-target mass
($\dtb$) stays small.
\textbf{(b)}~\emph{Off-target mass}: the mirror case, where off-target mass
($\dtb$) is large and coverage is perfect.
\textbf{(c)}~The two directed distances swap between cases and identify which
failure occurred, while a symmetric distance (their average) is nearly identical
for both. See Sect.~\ref{sec:setup} for the definitions.}
\label{fig:concept}
\end{figure}

\begin{table}[tbp]
\centering
\caption{\textbf{Transfer study.} Four architectures over 11 training
sources, evaluated by peak \pck{} at $\alpha{=}5\%$ on 9 target benchmarks. Backbones are frozen
(\encfrozen) or fine-tuned (\enctrained); RAFT has none and trains end-to-end.
\emph{Real-motion} datasets carry dense pixel-motion ground truth, \emph{semantic}
datasets use sparse keypoints.}
\label{tab:study}
\footnotesize
\centering
\setlength{\tabcolsep}{4pt}
% Two panels side by side. Nesting them in one outer tabular prevents the
% inter-panel glue from becoming a line break (which would stack them).
\begin{tabular}{@{}c@{\hspace{3mm}}c@{}}
\begin{tabular}[t]{@{}lll@{}}
\toprule
Architecture & Backbone & Encoder \\
\midrule
CATs++~\cite{cho2022cats++}            & ResNet-101 & \encfrozen\,/\,\enctrained \\
GLU-Net~\cite{truong2020glunet}        & VGG        & \encfrozen\,/\,\enctrained \\
FlowFormer~\cite{huang2022flowformer}  & Twins-SVT  & \encfrozen\,/\,\enctrained \\
RAFT~\cite{teed2020raft}               & none       & end-to-end \\
\bottomrule
\end{tabular}
&
\begin{tabular}[t]{@{}ll@{}}
\toprule
Real-motion & Semantic \\
\midrule
KITTI-12/15~\cite{kitti2012,kitti2015}            & SPair-71k~\cite{spair71k} \\
FlyingThings~\cite{mayer2016FlyingThingsDataset}  & PF-Pascal~\cite{pfpascal} \\
PointOdyssey~\cite{pointodyssey}                  & PF-Willow~\cite{pfpascal} \\
SDF-Fractals3D                                    & TSS~\cite{taniai2016tss} \\
\bottomrule
\end{tabular}
\end{tabular}

\smallskip
{\footnotesize Seven configurations $\times$ 11 training sources ($-$1: RAFT not trained on SPair) $=$
\textbf{76} trained models; \textbf{684} transfer measurements.}

\end{table}

In this work we make the following three contributions:
\begin{enumerate}
\item \textbf{A directional coverage rule} (Sect.~\ref{sec:law}): we find that source motion
  coverage of the target predicts transfer while appearance does not, and that ranking
  sources by this single directed distance which needs no training, nearly matches the
  ranking obtained by training a model on each source and measuring its transfer.
\item \textbf{From ranking to absolute prediction} (Sect.~\ref{sec:absolute}): we
  propose a few-shot adaptation of the zero-shot coverage rule that enables absolute
  \pck{} estimation.
\item \textbf{Coverage as a design objective} (Sect.~\ref{sec:interv}): we demonstrate
  that motion coverage can serve as a dataset-design objective, improving the transfer
  of two unrelated synthetic dataset generators, the photorealistic Kubric and an asset-free
  SDF-Fractals3D, to target tasks.
\end{enumerate}

